\chapter{Why translate Arabic logic? Burnett's 2004 transmission argument}
\label{sec:p13-burnett-arabic-latin-2004}

Latin scholars already possessed Boethian translations and explanations
of the older logical corpus, and could obtain further Aristotelian works
directly from Greek. Why, then, translate Arabic logical works whose
terminology and examples had been shaped for a different language?
Charles Burnett's 2004 chapter answers that question through the
curricula, interests and textual practices of particular translators.
This treatment follows that exact chapter, printed pp.~597--606,
not his broader companion survey of Arabic philosophy
\cite{Burnett2004ArabicLatinHHL}.

\section{The positive translation record and its limits}
\label{sec:p13-burnett-2004-translators}

Gerard of Cremona translated the \emph{Posterior Analytics},
Themistius' related commentary and an al-Farabi work called
\emph{De syllogismo}. The last work remained unidentified in
Latin in Burnett's account; his proposed association with
al-Farabi's commentary on demonstration is an inference, not
an identification of a newly located manuscript. The cluster
belongs to a Toledan programme of Arabic learning, not a
search limited to logical texts previously unavailable in Latin.
\cite[pp.~597--598, 604--605]{Burnett2004ArabicLatinHHL}

The chapter associates another Toledan project with Dominicus
Gundissalinus and Abraham ibn Daud, whose specimen translation
from Avicenna sought patronage. It associates Gundissalinus
and John of Spain with al-Ghazali's \emph{Aims of the Philosophers},
derived from Avicenna's Persian compendium. Avicenna's own
systematic writing, al-Ghazali's exposition and Averroes'
commentaries are thus distinct objects of translation.
\cite[p.~598]{Burnett2004ArabicLatinHHL}

A qualification is necessary here. The 2004 chapter's broad
description of the Avicennan collaboration includes the whole
logical division of the \emph{Shifa}. Burnett's separately
identified companion table differentiates particular logical
chapters, sample passages, excerpts and translators instead.
The broad phrase is not sufficient evidence of a complete
surviving medieval Latin version of Avicenna's modal syllogistic.
The two source descriptions must not be combined into that
stronger conclusion without the actual translation witnesses.
\cite[p.~598]{Burnett2004ArabicLatinHHL}
\cite[pp.~394--395]{Burnett2005ArabicLatinSupplement}

Averroes' middle commentaries on the \emph{Organon} included
Porphyry's \emph{Isagoge} and the \emph{Rhetoric} and \emph{Poetics}.
Burnett regards William of Luna's translations of the first three
logical commentaries at Naples in the 1220s as almost certain,
and his translations of the \emph{Prior} and \emph{Posterior
Analytics} as likely. A connection with Frederick II's new
university is inferred rather than established by a surviving
dedication. Siraj al-Din al-Urmawi reportedly wrote a logical
work at Frederick's court, but its identity is unknown. Hermann
the German's 1256 Toledan work and the possible identification
with Manfred's translator retain their separate evidential standing.
\cite[pp.~598--599]{Burnett2004ArabicLatinHHL}

Ramon Marti's citations and Ramon Llull's Catalan verse adaptation
show routes other than a complete literal Latin translation.
Marti's Galenic proof material occurs within a passage taken
from al-Razi's criticism of Galen; the nesting of authorities
matters. Jacob Anatoli translated logical middle commentaries
into Hebrew at Naples and acknowledged Frederick's patronage.
Later Hebrew-to-Latin work is therefore not the same route as
William's medieval Arabic-to-Latin work.
\cite[p.~599, notes~6--8]{Burnett2004ArabicLatinHHL}

\section{What the Giunta logical volumes actually contain}
\label{sec:p13-burnett-2004-giunta}

The Giunta edition of 1550--1552 exemplifies an aspiration to
assemble Aristotle and all available Averroist commentary.
Burnett's logical inventory contains six groups:
\cite[pp.~599--600]{Burnett2004ArabicLatinHHL}

\begin{enumerate}
\item The \emph{Isagoge}, \emph{Categories} and \emph{On
Interpretation}, with Averroes' middle commentary and
Gersonides' supercommentary, translated by Jacob Mantino.
\item The \emph{Prior Analytics} with a middle commentary
translated by Giovanni Francesco Burana.
\item The \emph{Posterior Analytics} with three Hebrew-mediated
long-commentary translations, by Abraham de Balmes, Burana
and Mantino. The parallel columns contract from three to two
where Mantino's translation stops. A middle commentary is also present.
\item Averroes' logical epitomes and questions, translated by de Balmes.
\item Letters and questions on particular logical subjects:
Averroes on primacy of demonstrative predicates; Abu al-Qasim
ibn Idris on genus and species; and texts under the names
Ibn Tumlus, Abu al-Qasim Muhammad ibn Qasim and an uncertain
Abu Abd al-Rahman ibn Jawhar. Their printed attribution is
reported, not converted into secure author identifications.
\item The \emph{Topics} and \emph{Sophistical Refutations}
with de Balmes' middle-commentary versions and an incomplete
Mantino \emph{Topics} version.
\end{enumerate}

Surviving manuscript counts do not alone measure readership.
Hissette's edition of the \emph{On Interpretation} commentary
identifies only three manuscripts yet traces use by Johannes
Pagus in 1229, Martin and John of Dacia, and an anonymous
writer. Burnett also records twelve printed Aristotle editions
containing the work between 1483 and 1560. Excerpts and
citations establish knowledge of other al-Farabi materials
without automatically establishing complete surviving versions.
The Brethren of Purity demonstration text's Latin attribution
to al-Kindi is another warning against reading author identity
straight off a received title.
\cite[pp.~600--601, notes~11--13]{Burnett2004ArabicLatinHHL}

\section{Literal translation can change the reader's problem}
\label{sec:p13-burnett-2004-calques}

Literalness was a deliberate policy for Gerard, Gundissalinus
and William, not merely an inability to write smooth Latin.
Burnett compares William's \emph{invenire}/\emph{inventum}
with the ordinary Latin vocabulary of being, and the use of
\emph{intentio} for Arabic \emph{ma'na}. A reader must decide
when a word concerns finding, existence, meaning or a subject
of predication. Mantino's alternative passage uses \emph{res}
and forms of \emph{esse}, bringing the exposition closer to
ordinary Latin expectations.
\cite[pp.~601--602, notes~14--15]{Burnett2004ArabicLatinHHL}

\begin{hbexample}{def:p13-burnett-2004-truth-bearer}{A name is not yet a truth-bearing assertion}
The compared passages distinguish simple names or verbal
expressions from a combination capable of truth or falsity.
A name for a human, a colour or a fabulous animal does not
assert that its referent exists merely by being uttered.
Adding an affirmation or denial of being, possibly with a
time qualification, changes the logical object. Reading
William's \emph{invenitur} here as a report of discovery
rather than as the intended expression of being would
change the proposition and its truth conditions. Mantino's
choice of \emph{est} avoids that particular lexical temptation;
it does not make every translation choice equivalent.
\cite[pp.~601--602]{Burnett2004ArabicLatinHHL}
\end{hbexample}

Burnett's other calques involve Greek terms rendered by
Arabic phrases and then translated literally: a hypothesis
becomes a placed root, \emph{radix posita}, and an enthymeme
a concealed syllogism, \emph{syllogismus occultus}. The occurrence
of \emph{syllogismus} in the latter phrase does not establish
that every property of a completed syllogism belongs to it.
The historical phrase is a route through technical vocabularies,
not a substitute for reconstructing the missing premise or
the intended inferential task.
\cite[p.~602, note~16]{Burnett2004ArabicLatinHHL}

The fabulous-animal example shows three distinct adaptations.
Averroes supplements Aristotle's goat-stag with an Arabic
fabulous creature. William transliterates that Arabic expression,
and Latin copying produces several corrupt forms. Burnett
suggests that William's particular form may reflect a variant
not attested in the Arabic manuscripts consulted. Mantino
instead uses the culturally familiar chimera. A replacement
can preserve the role of a name lacking an ordinary referent
without preserving the zoological identity of the example.
\cite[p.~602]{Burnett2004ArabicLatinHHL}

Arabic case morphology supplies a different failure. William's
reference to the case ending in \emph{u} reflects an Arabic
grammatical distinction, but copying turns the phrase into
unintelligible forms. Mantino glosses over that particular
ending. The literal token and the logical distinction between
an uninflected and an oblique expression are not the same
thing: omitting the opaque token avoids one obstruction but
does not establish a word-for-word equivalence.
\cite[p.~603, note~17]{Burnett2004ArabicLatinHHL}

\section{Audience adaptation and doctrinal alteration are different}
\label{sec:p13-burnett-2004-audience-doctrine}

Averroes' middle commentaries paraphrase Aristotle rather than
quote every sentence, aiming to clarify argumentative sequence
and make examples intelligible to an Arabic audience. William's
literal Latin versions show little evidence of equivalent
explanatory glossing. In the \emph{Poetics}, Averroes replaces
Greek poems with Arabic examples; Hermann translates those
Arabic examples rather than substitute Latin poetry. The
result carries the adapted Arabic illustrations into a new
audience instead of performing the adaptation anew.
\cite[pp.~603--604, notes~18--19]{Burnett2004ArabicLatinHHL}

Another example concerns doctrine, not merely style. Burnett
reports Hugonnard-Roche's comparison of the Greek and Arabic
\emph{Posterior Analytics} and Themistius: a paraphrastic route
changes the account of dialectic in relation to first principles.
Examining first principles is not the same task as demonstrating
the common propositions within each science. If a reader
adopts the second task from the translated wording, correcting
the Latin grammar alone cannot restore the first. This is a
specific reported textual comparison; it is not a claim that
every Arabic version makes the same alteration.
\cite[p.~604, note~20]{Burnett2004ArabicLatinHHL}

\section{The explanatory argument: curriculum and commentary}
\label{sec:p13-burnett-2004-explanation}

For Gerard, the relevant available scholarly culture at Toledo
was Arabic. Al-Farabi's \emph{Enumeration of the Sciences},
translated by both Gerard and Gundissalinus, presented logic
as preparation for mathematics, physics, metaphysics and
moral inquiry. Its logical chapter covered the prescribed
texts; translating demonstration and its commentary belonged
to replicating a coherent course of learning in Latin.
\cite[pp.~604--605]{Burnett2004ArabicLatinHHL}

Ahmad ibn Yusuf's work on ratio and proportion provides
a concrete connection to geometry. Its dialogue differentiates
practical mathematical competence and memorization from
understanding the principles of the art, for which it calls
on logic. Gerard translated the work and his associates used
it. Burnett accordingly connects the \emph{Posterior Analytics},
its Arabic title identifying demonstration, Themistius' help
with its difficulty, and the students' placement of logic
first in Gerard's bibliography. The bibliographical ordering
supports the curricular interpretation; it is not a dated
sequence of Gerard's translation activities.
\cite[pp.~604--605, notes~21--23]{Burnett2004ArabicLatinHHL}

William's motives remain conjectural because no dedication
or direct account of his activities supplies them. Burnett
suggests a wider programme to obtain Averroes' complete
exposition of Aristotle, together with Greek commentator
material otherwise only fragmentarily known. Association
of the logical texts with other Averroist commentaries,
Frederick's patronage and the later Giunta collection support
that explanation. They do not supply a recovered statement
of William's private intention.
\cite[pp.~605--606]{Burnett2004ArabicLatinHHL}

\begin{proofboundary}{def:p13-burnett-2004-boundary}{What the original chapter establishes}
The chapter's argument connects actual textual differences
and documented transmission to a qualified explanation
of translation. It supplies neither a universal translation
equivalence theorem nor complete editions of the compared
Arabic and Latin works. Its literal wording contains an
unresolved cross-reference in note~24; no page number is
invented here. Broad historical assertions, uncertain names
and proposed identifications retain their source-specific
standing. The recovered published source is distinct from
an independently verified history of every translation.
\end{proofboundary}
